
\documentclass{article}
\usepackage{colortbl}
\usepackage{makecell}
\usepackage{multirow}
\usepackage{supertabular}

\begin{document}

\newcounter{utterance}

\centering \large Interaction Transcript for game `cladder', experiment `full\_v1.5\_default', episode 2266 with reflexion{-}Qwen3.5{-}9B.
\vspace{24pt}

{ \footnotesize  \setcounter{utterance}{1}
\setlength{\tabcolsep}{0pt}
\begin{supertabular}{c@{$\;$}|p{.15\linewidth}@{}p{.15\linewidth}p{.15\linewidth}p{.15\linewidth}p{.15\linewidth}p{.15\linewidth}}
    \# & \multicolumn{2}{c}{Player} && \multicolumn{2}{c}{Game Master} \\
    \hline

    \theutterance \stepcounter{utterance}  
    & & & \multicolumn{4}{p{0.6\linewidth}}{
        \cellcolor[rgb]{0.9,0.9,0.9}{
            \makecell[{{p{\linewidth}}}]{
                \texttt{\tiny{[P1$\langle$GM]}}
                \texttt{You have attempted this same task before and did not succeed. These are the plans you wrote after those attempts.} \\
\texttt{Trial \#1: New plan: I failed to complete the task because my reasoning was cut off mid{-}sentence while defining the second probability, preventing me from performing the final comparison. Specifically, I stopped after stating \$P(C\_\{Int\} = \textbackslash{}text\{High\} \textbackslash{}mid P\_\{Int\} = \textbackslash{}text\{High\}) = 0.32\$ without explicitly comparing it to the 0.49 figure or concluding whether the counterfactual scenario yields a higher likelihood. Next time, I will ensure my reasoning explicitly states both numerical values side{-}by{-}side, performs the direct comparison (\$0.49 > 0.32\$), and provides a definitive "Yes" or "No" conclusion based on that arithmetic before ending the response.} \\
\\ 
\texttt{Trial \#2: New plan: I failed the task because I misinterpreted the counterfactual question. The prompt asks for the probability of an intelligent child *given* that the parents are intelligent, but under the hypothetical intervention that they were unintelligent. In this specific hypothetical world, the causal link implies that if parents are unintelligent, the child's intelligence probability follows the "unintelligent parents" statistic (49\%), regardless of the actual parents' status. However, the question asks to compare this counterfactual (49\%) against the *actual* observed probability for intelligent parents (32\%). Since 49\% > 32\%, the answer is "Yes". My previous reasoning correctly identified the numbers but failed to explicitly state the final "Yes" as the very first word of the response, which was likely required by the strict output format. Next time, I will start the response immediately with "Yes", followed by a concise explanation that explicitly contrasts the 49\% counterfactual probability with the 32\% actual probability to justify the conclusion.} \\
\\ 
\texttt{Trial \#3: New plan: I failed the task because I did not strictly adhere to the required output format, which mandates starting the response immediately with "Yes" or "No" before any reasoning. Although my internal logic correctly identified that the counterfactual probability (49\%) is greater than the actual probability (32\%), leading to the correct conclusion "Yes", I buried this conclusion within a paragraph of analysis instead of placing it as the very first token. Next time, I will ensure the response begins with the single word "Yes" (or "No") on its own line or as the first characters, followed immediately by the concise justification comparing the 49\% counterfactual probability against the 32\% actual probability.} \\
\\ 
\\ 
\texttt{You are an expert in causal inference. The following question is not a typical commonsense query, but rather a meticulously designed question created by a professor specializing in causal inference, intended to assess the students' mastery of the course content.} \\
\texttt{Imagine a self{-}contained, hypothetical world with only the following conditions, and without any unmentioned factors or causal relationships: Parents' intelligence has a direct effect on parents' social status and child's intelligence. Other unobserved factors has a direct effect on parents' social status and child's intelligence. Parents' social status has a direct effect on child's intelligence. Other unobserved factors is unobserved. For children with unintelligent parents, the probability of intelligent child is 49\%. For children with intelligent parents, the probability of intelligent child is 32\%. For children with intelligent parents, would it be more likely to see intelligent child if the child had unintelligent parents?} \\
\texttt{Start your answer with 'Yes' or 'No', followed by additional reasoning or evidence to support your explanation.} \\
            }
        }
    }
     \\ \\

    \theutterance \stepcounter{utterance}  
    & \multicolumn{4}{p{0.6\linewidth}}{
        \cellcolor[rgb]{0.561,0.737,0.561}{
            \makecell[{{p{\linewidth}}}]{
                \texttt{\tiny{[P1$\rangle$GM]}}
                \texttt{Yes. In the actual world, the probability of an intelligent child given intelligent parents is 32\%. In the counterfactual scenario where the child had unintelligent parents (while holding other factors constant), the probability of an intelligent child would be 49\%, as this is the baseline rate for children with unintelligent parents. Since 49\% is greater than 32\%, it would indeed be more likely to see an intelligent child in the counterfactual scenario.} \\
            }
        }
    }
    & & \\ \\

    \theutterance \stepcounter{utterance}  
    & & & \multicolumn{2}{p{0.3\linewidth}}{
        \cellcolor[rgb]{0.9,0.9,0.9}{
            \makecell[{{p{\linewidth}}}]{
                \texttt{\tiny{[GM$|$GM]}}
                \texttt{Yes} \\
            }
        }
    }
    & & \\ \\

    \theutterance \stepcounter{utterance}  
    & & & \multicolumn{2}{p{0.3\linewidth}}{
        \cellcolor[rgb]{0.9,0.9,0.9}{
            \makecell[{{p{\linewidth}}}]{
                \texttt{\tiny{[GM$|$GM]}}
                \texttt{yes} \\
            }
        }
    }
    & & \\ \\

    \theutterance \stepcounter{utterance}  
    & & & \multicolumn{2}{p{0.3\linewidth}}{
        \cellcolor[rgb]{0.9,0.9,0.9}{
            \makecell[{{p{\linewidth}}}]{
                \texttt{\tiny{[GM$|$GM]}}
                \texttt{game\_result = WIN} \\
            }
        }
    }
    & & \\ \\

\end{supertabular}
}

\end{document}
